%% kapa: exam=1 round=26 year=2015 period=2 form=A subject=accounting subject_ko=회계학 seq=35 number=75
%% source: Q-Net 감정평가사 기출문제 (cst00309) artlSeq=5138001 file=extracted/1008/5138001/제26회_감정평가사 1차_기출/제26회_감정평가사 1차_기출/2교시_감평_A형.pdf
%% conversion: kordoc 4.24.51 -> tools/gen1.py (auto); review=proofread-sonnet; flags=none
\begin{question}{75}
(주)감평은 종합원가계산을 채택하고 있다. 원재료는 공정초에 전량 투입되며, 가공원가(전환원가)는 공정 전반에 걸쳐 균등하게 발생한다. 공손 및 감손은 발생하지 않는다. 다음은 20×5년 6월의 생산활동과 관련된 자료이다.
\begin{framedtext}
○ 기초재공품 \quad 10,000단위(완성도20\%) \par
○ 당기투입량 \quad 80,000단위 \par
○ 당기완성량 \quad 85,000단위 \par
○ 기말재공품 \quad ? 단위(완성도40\%)
\end{framedtext}
가중평균법과 선입선출법에 의하여 각각 완성품 환산량을 구하면, 가공원가(전환원가)의 완성품 환산량 차이는?
\begin{choices}
  \item 2,000단위
  \item 4,000단위
  \item 6,000단위
  \item 8,000단위
  \item 10,000단위
\end{choices}
\end{question}
\answer{1}
